\documentclass[tikz,border=15pt]{standalone}
% Pacotes adicionados para garantir que os acentos compilem corretamente
\usepackage[utf8]{inputenc}
\usepackage[T1]{fontenc}
\usetikzlibrary{arrows.meta, positioning, shapes.geometric, calc}

\begin{document}

% =========================================================================
% PALETA DE CORES
% =========================================================================
\definecolor{CorPassoBorda}{HTML}{1A5276}       % Azul Escuro (Borda dos passos normais)
\definecolor{CorPassoFundo}{HTML}{EBF5FB}       % Azul Claro (Fundo dos passos normais)

\definecolor{CorDestaqueBorda}{HTML}{D35400}    % Laranja Escuro (Borda dos passos em destaque)
\definecolor{CorDestaqueFundo}{HTML}{FEF9E7}    % Laranja Claro (Fundo dos passos em destaque)

\definecolor{CorHubBorda}{HTML}{7D3C98}         % Roxo Escuro (Borda do "State of the Art")
\definecolor{CorHubFundo}{HTML}{F5EEF8}         % Roxo Claro (Fundo do "State of the Art")

\definecolor{CorSetaPrincipal}{HTML}{1A5276}    % Azul (Setas do fluxo principal)
\definecolor{CorSetaDestaque}{HTML}{D35400}     % Laranja (Setas de conexão com o Hub)
\definecolor{CorSetaFeedback}{HTML}{7D3C98}     % Roxo (Setas tracejadas de retorno)
% =========================================================================

\begin{tikzpicture}[
    node distance=0.8cm and 2.5cm,
    >=Stealth,
    thick,
    font=\sffamily\large,
    % Configuração dos Estilos
    stepbox/.style={rectangle, draw=CorPassoBorda, fill=CorPassoFundo, minimum width=4.2cm, minimum height=1cm, align=center, rounded corners=3pt},
    emphasis/.style={rectangle, draw=CorDestaqueBorda, fill=CorDestaqueFundo, minimum width=4.2cm, minimum height=1cm, align=center, rounded corners=3pt, font=\sffamily\bfseries\large},
    decisionbox/.style={diamond, aspect=2.2, draw=CorPassoBorda, fill=CorPassoFundo, minimum width=3.5cm, minimum height=1cm, align=center, inner sep=0pt},
    hub/.style={ellipse, draw=CorHubBorda, fill=CorHubFundo, minimum width=4.2cm, minimum height=1.2cm, align=center},
    % Estilos das Linhas/Setas
    mainline/.style={->, line width=1.2pt, draw=CorSetaPrincipal},
    highlightline/.style={->, line width=1.2pt, draw=CorSetaDestaque},
    feedbackline/.style={->, line width=1pt, draw=CorSetaFeedback, dashed}
]

    % --- COLUNA DA ESQUERDA: FLUXO PRINCIPAL ---
    \node[stepbox] (obs) {\textbf{Observation}};
    \node[emphasis, below=of obs] (search) {\textbf{Bibliography Research}};
    
    % (Corrigido) Sem "line breaks" dentro de \textbf{}
    \node[stepbox, below=of search] (hypothesis) {\textbf{Elaborate}\\ \textbf{Hypothesis}};
    \node[stepbox, below=of hypothesis] (experiment) {\textbf{Experiments}};
    \node[stepbox, below=of experiment] (analysis) {\textbf{Analyze Results}};
    
    \node[decisionbox, below=of analysis] (decision) {\textbf{Hypotheses}\\ \textbf{Hold}};
    
    \node[emphasis, below=of decision] (peer) {\textbf{Peer Review}};
    \node[emphasis, below=of peer] (publish) {\textbf{Publication}};

    % --- COLUNA DA DIREITA: CONTEXTO GLOBAL ---
    \node[hub, right=3.5cm of search] (literature) {\textbf{State of the Art}};
    \node[hub, above=1.5cm of literature] (researchers) {\textbf{Other Researchers}};

    % --- CONEXÕES E SETAS ---
    % Fluxo Principal
    \draw[mainline] (obs) -- (search);
    \draw[mainline] (search) -- (hypothesis);
    \draw[mainline] (hypothesis) -- (experiment);
    \draw[mainline] (experiment) -- (analysis);
    \draw[mainline] (analysis) -- (decision);
    
    % Saída SIM do losango
    \draw[mainline] (decision) -- node[right, font=\sffamily\small\bfseries] {\textbf{YES}} (peer);
    \draw[highlightline] (peer) -- (publish);

    % Interações com o "Estado da Arte"
    \draw[highlightline, <->] (search) -- node[above, font=\sffamily\small] {Analyze gaps} (literature);
    
    \draw[highlightline, <->] (researchers) -- node[right, font=\sffamily\small, align=center] {Analyze and\\ contribute} (literature);

    % Ligação Publicação -> Estado da Arte (Faz a curva 90 graus direitinho)
    \draw[highlightline] (publish.east) -| (literature.south);

    % Loop de Feedback (Saída NÃO do losango, contornando pela esquerda)
    \draw[feedbackline] (decision.west) -- ++(-1.8,0) node[above, font=\sffamily\small\bfseries] {\textbf{NO}\phantom{Nooo}} |- (hypothesis.west);

\end{tikzpicture}
\end{document}